% Part: first-order-logic 
% Chapter: axiomatic-deduction 
% Section: provability

% verification of properties of provability needed for maximally 
% consistent sets in the completeness chapter.

\documentclass[../../../include/open-logic-section]{subfiles}

\begin{document}

\iftag{FOL}
      {\olfileid{fol}{axd}{prv}}
      {\olfileid{pl}{axd}{prv}}

\olsection{أحكام \usetoken{S}{derivability}}

\begin{editorial}
نحمل في هذا القسم شرط الجدة في~\QR{} على الفروض المستعملة فعلًا في البرهان،
لا على ما لم يستعمل من أعضاء مجموعة محيطة. وهذا اصطلاح الدعم المنتهي الذي
تقوم عليه الرتابة والتراص.
\end{editorial}

\begin{prop}[الرتابة] متى كان~$\Gamma \subseteq \Delta$ و
$\Gamma \Proves !A$، كان كذلك~$\Delta \Proves !A$.
\end{prop}

\begin{proof}
كل~$\Gamma_0 \subseteq \Gamma$ منتهية داخلة أيضاً في~$\Delta$ دخول
مجموعة جزئية منتهية؛ فـ!!a{derivation} لـ$\Gamma_0 \Proves !A$ تشهد
كذلك بـ$\Delta \Proves !A$.
\end{proof}

\begin{prop}\ollabel{prop:provability}
يكون~$\Gamma \Proves !A$ إذا وفقط إذا كانت~$\Gamma\cup \{\lnot!A \}$
غير متسقة.
\end{prop}

\begin{proof}
أما أولاً فافرض~$\Gamma \Proves !A$. توجد إذن مجموعة منتهية
$\Gamma_0 \subseteq \Gamma$ تحقق~$\Gamma_0 \Proves !A$، فيلزم
$\Gamma_0 \cup \{\lnot !A\} \Proves \lfalse$، فتكون
$\Gamma \cup \{\lnot !A\}$ غير متسقة. وأما العكس، فلتكن
$\Gamma \cup \{\lnot !A\}$ غير متسقة؛ فعندئذ توجد مجموعة منتهية
$\Gamma_0 \subseteq \Gamma$ تحقق
$\Gamma_0 \cup \{\lnot !A\} \Proves \lfalse$. وتنتج مبرهنة الاستنباط
$\Gamma_0 \Proves \lnot !A \lif \lfalse$، وتعطينا \textbf{Ax0}
$\Gamma_0 \Proves (\lnot !A \lif \lfalse) \lif !A$. فنستخرج بـMP
$\Gamma_0 \Proves !A$، ثم بالرتابة~$\Gamma \Proves !A$.
\end{proof}

\begin{prop}\ollabel{prop:provability-properties}
\begin{enumerate}
\item \ollabel{prop:provability-contr} متى كان~$\Gamma \Proves !A$ و
$\Gamma \cup \{ !A\} \Proves \lfalse$، كانت~$\Gamma$ غير متسقة.

\item \ollabel{prop:provability-lnot} متى كان
$\Gamma \cup \{!A\} \Proves \lfalse$، كان
$\Gamma \Proves \lnot !A$.

\item \ollabel{prop:provability-exhaustive} متى كان
$\Gamma \cup \{!A\} \Proves \lfalse$ و
$\Gamma \cup \{\lnot !A\} \Proves \lfalse$، كان
$\Gamma \Proves \lfalse$.

\tagitem{prvOr}{\ollabel{prop:provability-lor-left} متى كان
$\Gamma \cup \{!A\} \Proves \lfalse$ و
$\Gamma \cup \{!B\} \Proves \lfalse$، كان
$\Gamma \cup \{!A \lor !B\} \Proves \lfalse$.}{}

\tagitem{prvOr}{\ollabel{prop:provability-lor-right} متى كان
$\Gamma \Proves !A$ أو~$\Gamma \Proves !B$، كان
$\Gamma \Proves !A \lor !B$.}{}

\tagitem{prvAnd}{\ollabel{prop:provability-land-left} متى كان
$\Gamma \Proves !A \land !B$، كان~$\Gamma \Proves !A$ و
$\Gamma \Proves !B$.}{}

\tagitem{prvAnd}{\ollabel{prop:provability-land-right} متى كان
$\Gamma \Proves !A$ و$\Gamma \Proves !B$، كان
$\Gamma \Proves !A \land !B$.}{}

\tagitem{prvIf}{\ollabel{prop:provability-mp} متى كان
$\Gamma \Proves !A$ و$\Gamma \Proves !A \lif !B$، كان
$\Gamma \Proves !B$.}{}

\tagitem{prvIf}{\ollabel{prop:provability-lif} متى كان
$\Gamma \Proves \lnot !A$ أو~$\Gamma \Proves !B$، كان
$\Gamma \Proves !A \lif !B$.}{}
\end{enumerate}
\end{prop}

\begin{proof}
\begin{enumerate}
\item عيّن مجموعتين منتهيتين~$\Gamma_0,\Gamma_1 \subseteq \Gamma$
تحققان~$\Gamma_0 \cup \{!A\} \Proves \lfalse$ و
$\Gamma_1 \Proves !A$.

\begin{derivation}
1. & $\Gamma_0 \cup \{!A\} \Proves \lfalse$ & فرض \\
2. & $\Gamma_1 \Proves !A$ & فرض \\
3. & $\Gamma_0 \cup \Gamma_1 \Proves \lfalse$ & القطع \\
\end{derivation}

وبما أن~$\Gamma_0 \subseteq \Gamma$ و$\Gamma_1 \subseteq \Gamma$، فإن
$\Gamma_0 \cup \Gamma_1 \subseteq \Gamma$، فينتج
$\Gamma \Proves \lfalse$.

\item عيّن مجموعة منتهية~$\Gamma_0 \subseteq \Gamma$ تحقق
$\Gamma_0 \cup\{!A\} \Proves \lfalse$.

\begin{derivation}
1. & $\Gamma_0 \cup\{!A\} \Proves \lfalse$ & فرض \\
2. & $\Gamma_0 \Proves !A \lif \lfalse$ & مبرهنة الاستنباط \\
3. & $\Gamma_0 \Proves (!A \lif \lfalse) \lif \lnot !A$ & Ax0 \\
4. & $\Gamma_0 \Proves \lnot !A$ & MP 2, 3 \\
\end{derivation}

\item عيّن مجموعتين منتهيتين~$\Gamma_0,\Gamma_1 \subseteq \Gamma$
تحققان~$\Gamma_0 \cup \{ !A \} \Proves \lfalse$ و
$\Gamma_1 \cup \{\lnot !A \} \Proves \lfalse$.

\begin{derivation}
1. & $\Gamma_0 \cup\{!A\} \Proves \lfalse$ & فرض \\
2. & $\Gamma_0 \Proves !A \lif \lfalse$ & مبرهنة الاستنباط \\
3. & $\Gamma_1 \cup \{\lnot !A \} \Proves \lfalse$ & فرض \\
4. & $\Gamma_0 \Proves (!A \lif \lfalse) \lif \lnot !A$ & Ax0 \\
5. & $\Gamma_0 \Proves \lnot !A$ & MP 2, 4 \\
6. & $\Gamma_0 \cup \Gamma_1 \Proves \lfalse$ & القطع 3، 5 \\
\end{derivation}

ولما كان~$\Gamma_0 \subseteq \Gamma$ و$\Gamma_1 \subseteq \Gamma$، كان
$\Gamma_0 \cup \Gamma_1 \subseteq \Gamma$؛ ولذلك
$\Gamma \Proves \lfalse$.

% prop:provability-lor-left 
\tagitem{defOr}{}{
\iftag{probOr}{تمرين.}{تمرين.}}

% prop:provability-lor-right 
\tagitem{defOr}{}{\iftag{probOr}{تمرين.}{
لنفرض~$\Gamma_0 \Proves !A$.

\begin{derivation}
1. & $\Gamma_0 \Proves !A$ & فرض \\
2. & $\Gamma_0 \Proves !A \lif (!A \lor !B)$ & Ax0 \\
3. & $\Gamma_0 \Proves !A \lor !B$ & MP 1, 2 \\
\end{derivation}

فينتج~$\Gamma \Proves !A \lor !B$. ونظير ذلك يبرهن حالة
$\Gamma \Proves !B$.}}

% prop:provability-land-left 
\tagitem{defAnd}{}{\iftag{probAnd}{تمرين}{
لنفرض~$\Gamma_0 \Proves !A \land !B$.

\begin{derivation}
1. & $\Gamma_0 \Proves !A \land !B$ & فرض \\
2. & $\Gamma_0 \Proves (!A\land !B) \lif !A $ & Ax0 \\
3. & $\Gamma_0 \Proves !A$ & MP 1, 2 \\
\end{derivation}

فيلزم~$\Gamma \Proves !A$. وعلى مثاله يبيّن~!!{derivation} آخر أن
$\Gamma \Proves !B$.}}

% prop:provability-land-right
\tagitem{defAnd}{}{\iftag{probAnd}{تمرين.}{تمرين.}}

% prop:provability-mp 
\tagitem{defIf}{}{\iftag{probIf}{تمرين.}{لنفرض
$\Gamma_0 \Proves !A$ و$\Gamma_0 \Proves !A \lif !B$ معًا.

\begin{derivation}
1. & $\Gamma_0 \Proves !A$ & فرض \\
2. & $\Gamma_0 \Proves !A \lif !B $ & فرض \\
3. & $\Gamma_0 \Proves !B$ & MP 1, 2 \\
\end{derivation}

ولأن~$\Gamma_0 \subseteq \Gamma$، لزم~$\Gamma \Proves !B$.}}

% prop:provability-lif 
\tagitem{defIf}{}{\iftag{probIf}{تمرين.}{افرض أولاً
$\Gamma \Proves \lnot !A$.

\begin{derivation}
1. & $\Gamma_0 \Proves \lnot !A$ & فرض \\
2. & $\Gamma_0 \Proves \lnot !A \lif (!A \lif !B) $ & Ax0 \\
3. & $\Gamma_0 \Proves !A \lif !B$ & MP 1, 2 \\
\end{derivation}

وأما حالة~$\Gamma \Proves !B$ فبرهانها نظير ذلك:

\begin{derivation}
1. & $\Gamma_0 \Proves !B$ & فرض \\
2. & $\Gamma_0 \Proves !B \lif (!A \lif !B) $ & Ax0 \\
3. & $\Gamma_0 \Proves !A \lif !B$ & MP 1, 2 \\
\end{derivation}}}

\end{enumerate}
\end{proof}

\begin{probtag}{probOr,probAnd,probIf}
أتمم برهان~\olref[fol][axd][prv]{prop:provability-properties}.
\end{probtag}

\begin{thm}[التعميم]\ollabel{thm:Generalization}
متى كان~$\Gamma \Proves !A$، ولم يكن~$x$ حراً في شيء من~!!{formula}
$\Gamma$، كان~$\Gamma \Proves \lforall[x][!A]$.
\end{thm}

\begin{proof}
خذ برهانًا منتهيًا لـ~$!A$ من~$\Gamma$، ولتكن~$\Gamma_0$ مجموعة
!!{formula}s~$\Gamma$ التي عمل بها. وعيّن~!!a{constant}~$c\in\mathcal E$ لا يقع في
شيء من صيغ البرهان، ولا يساوي أي ثابت مميز عمل لتسويغ تطبيق لـ~\QR{} فيه،
ثم أحل~$c$ محل كل وقوع حر لـ~$x$ في كل خطوة. ولما كان
$x$ غير حر في~$\Gamma_0$ بقيت خطوات الفرض على حالها. ويظهر بالاستقراء في
خطوات البرهان أن المتتالية المستحدثة برهان لـ
$\Subst{!A}{c}{x}$ من~$\Gamma_0$: فالتعويض يحفظ أمثلة مسلمات الروابط
والمكمّمات ومسلّمات الهوية ذات الحدود المغلقة، ويحفظ MP وصورتي~\QR{}،
ويحفظ كذلك شرط الجدة فيهما لأن~$c$ اختير جديدًا. فيكون
$\Gamma_0 \Proves \Subst{!A}{c}{x}$.

ثم نشتق:
\begin{derivation}
1. & $\Gamma_0 \Proves \Subst{!A}{c}{x}$ & كما تقدم \\
2. & $\Gamma_0 \Proves \Subst{!A}{c}{x} \lif
  (\ltrue \lif \Subst{!A}{c}{x})$ & \olref[prp]{ax:lif1} \\
3. & $\Gamma_0 \Proves \ltrue \lif \Subst{!A}{c}{x}$ & MP 1، 2 \\
4. & $\Gamma_0 \Proves \ltrue \lif \lforall[x][!A]$ & \QR{} 3 \\
5. & $\Gamma_0 \Proves \ltrue$ & \olref[prp]{ax:ltrue} \\
6. & $\Gamma_0 \Proves \lforall[x][!A]$ & MP 4، 5
\end{derivation}
وإنما جاز~\QR{} في الخطوة~4 لأن~$c$ لا يقع في~$\Gamma_0$ ولا في~$\ltrue$
ولا في~$!A$.
ثم تعطينا الرتابة~$\Gamma \Proves \lforall[x][!A]$، وهو المطلوب.
\end{proof}

\begin{thm}\ollabel{thm:WeakGen}
(\emph{التعميم الضعيف على الثوابت}) إذا كان~$\Gamma \Proves !A$، وكان
$c$ !!a{constant} لا يقع في~$\Gamma$، وُجد متغير~$x$ لا يقع في~$!A$
يحقق~$\Gamma \Proves \lforall[x][\Subst{!A}{x}{c}]$، ويخلو البرهان
من~$c$.
\end{thm}

\begin{proof}
% تصحيح برهاني (0643-I2): عولجت مسلّمتا المكمّمات وصورتا QR بتجديد الثوابت ثم تبديل ثابتين تقابليًا.
خذ برهانًا منتهيًا~$\Pi$ لـ~$!A$ من~$\Gamma$، وليكن~$\Gamma_0$
دعمه المتناهي. طبّق \olref[qua]{lem:qr-freshening} على~$\Pi$ مع طائفة
المنع~$\{c\}$؛ فينتج برهان~$\Pi^*$ خاتمته~$!A$ ودعمه داخل
في~$\Gamma_0$، ولا يساوي~$c$ شيء من ثوابته المميزة. ولما كان~$\Pi^*$
متناهيًا، اختر~$d\in\mathcal E$ لا يقع في شيء من صيغه ولا يساوي ثابتًا
مميزًا على عقدة~\QR{} فيه. وعيّن متغيرًا~$x$ لا يقع في~$!A$.

ولتكن~$\rho$ إعادة التسمية المنتظمة التقابلية للثوابت التي تبدّل~$c$ و~$d$
وتثبت ما سواهما؛ وضع~$!A_d=\rho(!A)$، واكتب هذا اختصارًا~$!A[d/c]$
لأن~$d$ لم يقع في~$!A$. فإذا أجرينا~$\rho$ في كل سطر
من~$\Pi^*$ حصلنا على برهان لـ~$\rho(!A)$ من~$\Gamma_0$. فدعم هذا البرهان باق لأن~$c$
لا يقع في~$\Gamma$، ومسلمات الروابط محفوظة. وتبقى مسلّمتا المكمّمات حالتين
منهما، إذ تبعث~$\rho$ كل حد مغلق~$t$ إلى الحد المغلق~$\rho(t)$؛ وتحفظ
كذلك مسلمات الهوية ذات الحدود المغلقة، وتحفظ~MP. وأما تطبيق~\QR{} ذو
الثابت المميز~$a$، فإن~$a\ne c,d$ بالتجديد؛ فتحفظ~$\rho$ مقدمته
وخاتمته وشرط الجدة في كلتا صورتي~\QR{}.
ويخلو البرهان الناتج من~$c$.

ضع~$!A_x=\Subst{!A}{x}{c}$. ولأن~$x$ و~$d$ لا يقعان في~$!A$، فإحلال~$d$
محل~$x$ في~$!A_x$ هو~$!A_d=\rho(!A)=!A[d/c]$ نفسه. ومن البرهان المتقدم نشتق:
\begin{derivation}
1. & $\Gamma_0 \Proves !A_d$ & كما تقدم \\
2. & $\Gamma_0 \Proves !A_d \lif (\ltrue \lif !A_d)$
   & \olref[prp]{ax:lif1} \\
3. & $\Gamma_0 \Proves \ltrue \lif !A_d$ & MP 1، 2 \\
4. & $\Gamma_0 \Proves \ltrue \lif \lforall[x][!A_x]$
   & \QR{} 3، بالثابت المميز~$d$ \\
5. & $\Gamma_0 \Proves \ltrue$ & \olref[prp]{ax:ltrue} \\
6. & $\Gamma_0 \Proves \lforall[x][!A_x]$ & MP 4، 5
\end{derivation}
ويجوز~\QR{} لأن~$d$ لا يقع في~$\Gamma_0$ ولا في~$\ltrue$ ولا في~$!A_x$.
ثم تعطينا الرتابة
$\Gamma \Proves \lforall[x][\Subst{!A}{x}{c}]$ ببرهان يخلو من~$c$.
\end{proof}

\begin{explain}
نريد أن نستبدل بشرط~\olref{thm:WeakGen}، وهو ألا يقع~$x$ في~$!A$
البتة، شرطين أضعف منه: ألا يكون~$x$ حراً في~$!A$، وأن يكون~!!{free for}
$c$ فيها. ويتأتى ذلك بتبديل~!!{variable} مقيد، فنثبت هذا أولاً.
\end{explain}

\begin{lem}
\ollabel{lem:free-for}
متى كان~$x$ حراً لـ$c$ في~$!A$، ولم يكن~$y$ حراً في~$!A$، كان~$x$
!!{free for}~$y$ في~$\Subst{!A}{y}{c}$.
\end{lem}

\begin{proof}
لو لم يكن~$x$ !!{free for}~$y$ في~$\Subst{!A}{y}{c}$، لوقع ظهور حر
لـ$y$ في~$\Subst{!A}{y}{c}$ داخل مجال السور~$\lforall[x]$. ولكن~$y$
ليس حراً في~$!A$ بحكم الفرض، فكل ظهور من هذه الظهورات آت من التعويض
$\Subst{}{y}{c}$. فيقع إذن ظهور لـ$c$ داخل مجال~$\lforall[x]$، ولا
يكون~$x$ !!{free for}~$c$ في~$!A$.
\end{proof}

\begin{lem}[تغيير المتغير المقيد]
\ollabel{lem:ChangeBdVar}
متى لم يكن~$x$ ولا~$y$ حراً في~$!A$، وكان كل منهما~!!{free for}~$c$ في
$!A$، كان
$\Proves \lforall[x][\Subst{!A}{x}{c}] \equiv
\lforall[y][\Subst{!A}{y}{c}]$، ولا يتضمن البرهان~$c$.
\end{lem}

\begin{proof}
بحسب الفروض، يكون~$y$ !!{free for}~$c$ في~$!A$، ولا يكون~$x$ حراً
في~$!A$؛ فباللمّة المتقدمة~\olref{lem:free-for} يكون~$y$
!!{free for}~$x$ في~$\Subst{!A}{x}{c}$. فتكون
$\lforall[x][\Subst{!A}{x}{c} \lif
\Subst{!A}{x}{c} \Subst{}{y}{x}]$ مسلّمة (\textbf{Ax1}). ولما كان~$x$
غير حر في~$!A$، كان
$\Subst{!A}{x}{c} \Subst{}{y}{x} = \Subst{!A}{y}{c}$؛ ومن ثم
\[
\Proves \lforall[x][\Subst{!A}{x}{c}] \lif \Subst{!A}{y}{c}.
\]
ثم تعطينا مبرهنة الاستنباط
$\lforall[x][\Subst{!A}{x}{c}] \Proves \Subst{!A}{y}{c}$. ولما كان~$y$
غير حر في~$!A$، لم يكن حراً كذلك في
$\lforall[x][\Subst{!A}{x}{c}]$؛ فيعطينا التعميم، مع اختيار ثابته
المميز المستجد من~$\mathcal E\setminus\{c\}$،
$\lforall[x][\Subst{!A}{x}{c}] \Proves
\lforall[y][\Subst{!A}{y}{c}]$، ثم تعود مبرهنة الاستنباط فتعطي
$\Proves \lforall[x][\Subst{!A}{x}{c}] \lif
\lforall[y][\Subst{!A}{y}{c}]$. وأما الجهة الأخرى
$\Proves \lforall[y][\Subst{!A}{y}{c}] \lif
\lforall[x][\Subst{!A}{x}{c}]$ فبرهانها متناظر تماماً، مع اختيار الثابت
المميز على الوجه نفسه. وجميع الصيغ والخطوات المتقدمة خالية من~$c$، كما أن
بناء مبرهنة الاستنباط لا يستحدثه؛ فتنتج القضية بـTaut ببرهان يخلو من~$c$.
\end{proof}

\begin{thm}[التعميم القوي على الثوابت]
\ollabel{thm:StrongGen}
متى كان~$\Gamma \Proves !A$، ولم يقع~!!{constant}~$c$ في~$\Gamma$،
وكان~$x$ غير حر في~$!A$ ولكنه~!!{free for}~$c$ فيها، كان
$\Gamma \Proves \lforall[x][\Subst{!A}{x}{c}]$.
\end{thm}

\begin{proof}
لما كان~$\Gamma \Proves !A$، وكان~$c$ غير واقع في~$\Gamma$، أوجد لنا
التعميم الضعيف~!!a{variable}~$y$ لا يقع في~$!A$، بحيث
$\Gamma \Proves \lforall[y][\Subst{!A}{y}{c}]$. ولأن~$y$ لا يقع في
$!A$ البتة، فهو غير حر فيها، وهو~!!{free for}~$c$ فيها أيضاً. وإذا كان
$x$، على ما في فرض المبرهنة، غير حر في~$!A$ و!!{free for}~$c$ فيها،
تمت شروط تبديل~!!{variable} مقيد، فيكون
$\Proves \lforall[x][\Subst{!A}{x}{c}] \equiv
\lforall[y][\Subst{!A}{y}{c}]$، ويلزم منه
$\Gamma \Proves \lforall[x][\Subst{!A}{x}{c}]$.
\end{proof}

\begin{thm}
\ollabel{thm:provability-quantifiers}
\begin{tagenumerate}{prvEx,prvAll}
\tagitem{prvEx}{متى كان~$\Gamma \Proves !A(t)$، كان
$\Gamma \Proves \lexists[x][!A(x)]$.}{}
\tagitem{prvAll}{متى كان~$\Gamma \Proves \lforall[x][!A(x)]$، كان
$\Gamma \Proves !A(t)$.}{}
\end{tagenumerate}
\end{thm}

\begin{proof}
\begin{tagenumerate}{prvEx,prvAll}
%%Need axioms for exists
\tagitem{prvEx}{تمرين.}{}
\tagitem{prvAll}{لنفرض~$\Gamma_0 \Proves \lforall[x][!A(x)]$.

\begin{derivation}
1. & $\Gamma_0 \Proves \lforall[x][!A(x)]$ & فرض \\
2. & $\Gamma_0 \Proves \lforall[x][!A(x)] \lif \Subst{!A}{t}{x}$ & Ax1 \\
3. & $\Gamma_0 \Proves \Subst{!A}{t}{x}$ & MP 1, 2 \\
\end{derivation}}{}

\end{tagenumerate}
\end{proof}

\end{document}
